% Privacy-first section 4: released models leak at a quality a practitioner would accept.
% Budget 1.75 pages including Table~\ref{tab:cost} and Figure~\ref{fig:cost}.
% Every number here is from the single-basis re-measurement of 19-23 Oct; see notes 3.58-3.59.
\section{Released models leak at a quality you would ship}
\label{sec:leak}

Whether a model memorises after being trained past the point of usefulness is not the question a
practitioner faces. The question is whether leakage appears at a checkpoint they would choose. We
therefore report every leak twice: as a rate, and as the held-out loss it costs relative to the
best checkpoint that run itself produced. The second number is what makes the first actionable,
because a leak that costs a hundredth of a bit is one nobody would trade away.

\paragraph{The cheapest leak costs seven thousandths of a bit.}
Table~\ref{tab:cost} orders the released checkpoints by that price. DNABERT-2 reproduces 30\% of
planted records at \textbf{+0.007\,\bits} above its own optimum. No one selecting a checkpoint
would notice that difference, and none of the six downstream tasks we measure moves by more than one per
cent across the same interval (\S\ref{sec:blind}). Nucleotide Transformer gives up
\textbf{+0.03\,\bits} to leak 8--14\% across three seeds. Both are far inside the range a user
would consider interchangeable with the optimum.

The control is flat throughout. Records the model never saw are reproduced 0\% of the time in every
cell, every tier and every seed, and the never-trained window at the foot of each block is scored by
the same decoder on the same host, so the zero is a property of the measurement and not of a
separate baseline.

\begin{table}[t]
\centering
\caption{Released masked checkpoints on \emph{E.~coli}, records planted $16\times$, 50 probes per
tier, three seeds where shown. \emph{Cost} is held-out \bits{} above that run's own best saved
checkpoint. \emph{Standard} is the prefix-conditioned protocol the field uses; \emph{ours} is the
re-masking decoder of \S\ref{sec:measure}. All cells re-measured on one numerical basis.}
\label{tab:cost}
\begin{tabular}{llccc}
\toprule
checkpoint & corpus & cost & standard & ours \\
\midrule
DNABERT-2 117M      & \emph{E.~coli} & \textbf{+0.007} & \textbf{30\%}      & 0\% \\
NT 250M (3 seeds)   & \emph{E.~coli} & \textbf{+0.03}  & 4 / 8 / 4\%        & \textbf{8 / 8 / 14\%} \\
NT 250M (3 seeds)   & \emph{E.~coli} & +0.44           & 26 / 26 / 30\%     & \textbf{66 / 68 / 76\%} \\
NT 500M             & \emph{E.~coli} & +0.46           & 36\%               & \textbf{78\%} \\
\midrule
never-trained control & --- & --- & 0\% & 0\% \\
\bottomrule
\end{tabular}
\end{table}

\paragraph{The two protocols disagree, and neither dominates.}
Reading across Table~\ref{tab:cost} shows why a single extraction test cannot carry this
measurement. On fixed-width non-overlapping tokenisers the standard protocol under-reports by
roughly a factor of two: at +0.44\,\bits{} it recovers 26--30\% where ours recovers 66--76\%. On
DNABERT-2 the order reverses completely --- the standard protocol recovers 30\% and ours recovers
nothing at all. We return to this in \S\ref{sec:measure}; here it matters only that the figures in
this table are each read with whichever protocol is stronger for that model, so they are lower
bounds on what an attacker could obtain and not artefacts of one decoder's blind spot.

\paragraph{On real human haplotypes a single copy is enough.}
The \emph{E.~coli} canaries above are uniform random sequence, which makes them a clean lower bound
and a fair target for criticism: a random string shares nothing with the statistics of a genome. We
therefore repeat the measurement with canaries that are real. Each is a 96\nt{} window of human
chromosome~22 built around a variant carried by at most five of the 2{,}548 individuals in
1000~Genomes phase~3, written with that carrier's own alleles. Reproducing one is reproducing a
specific person's rare haplotype, not a string.

\begin{table}[t]
\centering
\caption{Nucleotide Transformer fine-tuned on human chromosome~22 with real rare-variant haplotypes
as canaries. \emph{1 copy} is the realistic case: one individual's haplotype appearing once in the
fine-tuning set.}
\label{tab:human}
\begin{tabular}{lccccc}
\toprule
& & \multicolumn{2}{c}{16 copies} & \multicolumn{2}{c}{1 copy} \\
\cmidrule(lr){3-4}\cmidrule(lr){5-6}
checkpoint & held-out \bits & standard & ours & standard & ours \\
\midrule
NT 250M, round 15 & 1.145 & 6\%  & 6\%           & 0\% & 0\% \\
NT 250M, round 30 & 0.739 & 10\% & 28\%          & 0\% & 0\% \\
NT 250M, round 50 & 0.465 & 18\% & \textbf{52\%} & 2\% & \textbf{8\%} \\
NT 500M, round 15 & 1.045 & 16\% & 14\%          & 0\% & 0\% \\
NT 500M, round 50 & 0.396 & 24\% & \textbf{58\%} & 0\% & \textbf{12\%} \\
\midrule
never-trained control & --- & 0\% & 0\% & 0\% & 0\% \\
\bottomrule
\end{tabular}
\end{table}

Table~\ref{tab:human} is the result we would ask a reader to carry away from this section. A
haplotype belonging to one person, present exactly once in the fine-tuning data, is returned letter
for letter 8\% of the time by the 250M model and 12\% by the 500M --- roughly one case in ten, on
weights anyone can download, with the matched never-trained control at zero. The repeated tier
behaves as the \emph{E.~coli} arm does, so the single-copy column is not an artefact of a canary
appearing often enough to be learned as vocabulary.

\paragraph{What the model knows is not what our decoder can reach.}
One column of Table~\ref{tab:cost} deserves a caution, because it is the one most easily misread.
Our decoder recovers 0\% from DNABERT-2, and that is not evidence the model holds nothing. Asked for
any single nucleotide of a planted record with every other nucleotide supplied, DNABERT-2 answers
correctly 93--95\% of the time for records planted \emph{once}, which puts the implied ceiling on
whole-record recovery at 0.21--0.39 across three seeds. The information is demonstrably present; our
search does not find it, because a variable-length tokenisation re-segments the sequence on every
proposal and leaves the search no foothold. The standard protocol, which is matched to that vocabulary,
reaches 30\%. We report the larger of the two, and note that taking the smaller one for the model's
exposure is the mistake this paper is about.
